% !TeX root = main.tex
\lecture{1}{Tue 20 Jan 2026 12:00}{Course Welcome and Intro to Matter}

\section{Course Welcome}
\begin{itemize}
    \item Study guide and full revision notes (from Prof. Mike Gunn) on Canvas.
    \begin{itemize}
        \item The course has changed a bit recently, so the full revision notes have additional theory/content which isn't examinable.
    \end{itemize}
    \item Series of non-assessed problems also on Canvas.
\end{itemize}

\subsection{Course Aims}
\begin{tcolorbox}
    We want to be able to explain the properties of matter (in the macroscopic world) of various kinds, based on their microscopic behaviour. We also want to be able to understand more about the microscopic behaviour of matter from macroscopic observables.
\end{tcolorbox}

\subsection{Recommended Reading}
\begin{itemize}
    \item Young and Freedman, Ch14 and Ch17-19.
    \item Properties of Matter, Flowers and Mendoza.
    \item Introduction to Thermal Physics, Adkins.
\end{itemize}

\section{Intro to Matter}
\subsection{States of Matter}
We want to move away from a simple "cold stuff is a solid, warmer stuff is a liquid and hotter stuff is a gas" by also considering how different materials change state as a function of both pressure and temperature:
\begin{figure}[H]
    \centering
    \includegraphics[width=0.75\textwidth]{figures/lec01-01.png}
     \caption{Phase Diagram for Water}
\end{figure}
\begin{figure}[H]
    \centering
    \includegraphics[width=0.75\textwidth]{figures/lec01-02.png}
     \caption{Phase Diagram for CO2}
\end{figure}

At ``Standard Pressure and Temperature'' (at $20 \si{\degreeCelsius}$ and $1$ATM), we see that increasing pressure of C02 takes it to a liquid state. However, at lower temperatures, decreasing the pressure can cause it to sublimate straight from a solid to a gas, without a liquid intermediary.

Taking water as an example, pushing the pressure/temperature causes another high-energy state altogether called a ``supercritical fluid''. The gradient for the initial vertical portion of the solid/liquid boundary around the freezing point is slightly negative, so applying pressure to ice on this boundary causes it to become liquid. This is uncommon for many materials. 

This is because water is weird, where the solid form is less dense than the liquid form. This is how ice skating works, where applying pressure (from a skate) causes a thin upper layer of ice to melt. This thin water layer is what the ice skate can travel on.

Every material also has a ``triple point'' where the material can act simultaneously in all three standard states of matter. 

